\documentclass[10pt,a4paper]{article}
\usepackage[T1]{fontenc}
\usepackage{lmodern}
\usepackage[margin=19mm]{geometry}
\usepackage{amsmath,amssymb,booktabs,tabularx,array}
\usepackage[ruled,vlined,linesnumbered]{algorithm2e}
\usepackage[hidelinks]{hyperref}
\setlength{\parindent}{0pt}
\setlength{\parskip}{4pt}
\setlength{\tabcolsep}{5pt}
\SetKwInOut{Input}{Input}
\SetKwInOut{Output}{Output}
\SetKw{Return}{return}
\SetKw{Continue}{continue}
\SetKw{Break}{break}
\SetAlFnt{\small}
\SetAlgoNlRelativeSize{-1}
\SetAlCapFnt{\bfseries}
\SetAlCapNameFnt{\bfseries}
\newcommand{\Pop}{\mathcal P}
\newcommand{\Reg}{\mathcal Q}
\newcommand{\Hist}{\mathcal H}
\newcommand{\Trace}{\mathcal Z}
\newcommand{\LLM}{\mathcal L_{\phi}}
\newcommand{\Native}{\operatorname{Native}}
\newcommand{\Best}{\operatorname{best}}
\newcommand{\Adapt}{\pi_{\mathrm{adapt}}}
\newcommand{\Blank}{\varnothing}
\newcommand{\Ind}{\mathbf 1}

\begin{document}
\begin{center}
{\Large\bfseries CIPHEUR C05: Problem Model and Algorithms}\\[3pt]
{\small Fixed-contact scheduling; finite structural plans; asynchronous selection}
\end{center}

\textbf{Problem model.}
An instance is $\mathcal I=(\mathcal S,\mathcal A,V,\delta^A,\delta^S)$,
where $\mathcal S$ and $\mathcal A$ are satellites and antennas. Each contact
$v=(s_v,a_v,b_v,e_v)\in V$ is a complete predefined visibility window, with
fixed satellite, antenna, start time, and end time. The decision is to select
or omit the entire contact; its start and duration are not decision variables.
Its positive weight is $w_v=e_v-b_v$, represented in integer time ticks.
The resource separation requirements are $\delta^A_a$ and $\delta^S_s$.

Order two distinct contacts as $i\prec j$ by $(b_i,i)<(b_j,j)$ and define
$\Delta_{ij}=b_j-e_i$. The conflict graph is $G=(V,E)$, where
\begin{align*}
E_A&=\{\{i,j\}:i\prec j,\ a_i=a_j,\ \Delta_{ij}<\delta^A_{a_i}\},\\
E_S&=\{\{i,j\}:i\prec j,\ s_i=s_j,\ \Delta_{ij}<\delta^S_{s_i}\},
\qquad E=E_A\cup E_S.
\end{align*}
Equality to one resource's required gap is compatible for that resource;
the other shared resource can still introduce a conflict. The supplied
instance has no separate station or pass-identity conflict rule.
The scheduling problem is
\begin{equation}
\max_{x\in\{0,1\}^{|V|}}\sum_{v\in V}w_vx_v
\quad\text{subject to}\quad x_i+x_j\leq1\quad(\{i,j\}\in E).
\label{eq:model}
\end{equation}
Equivalently, select an independent set $X\subseteq V$ and maximize
$F(X)=\sum_{v\in X}w_v$. Graph edges and original weights remain the
authority for every feasibility check and objective evaluation.

\textbf{Unified notation.}
\begin{center}\small
\begin{tabularx}{\linewidth}{@{}p{.25\linewidth}X@{}}
\toprule
Symbol & Meaning \\
\midrule
$G=(V,E)$; $X$; $F(X)$ & Conflict graph; feasible schedule; original integer objective.\\
$\Pop=(X_0,\ldots,X_{m-1})$ & Persistent population of feasible search trajectories.\\
$X^*$; $j_{\Best}(\Pop)$ & Historical incumbent; current best population slot (lowest index breaks ties).\\
$\Trace$; $\Hist$; $\Hist_k$ & Native progress trace; closed interval ledger; its bounded visible history at opportunity $k$.\\
$o_k$; $\Reg_k$ & Observation and finite plan registry bound to decision opportunity $k$.\\
$\pi$; $\pi_k$; $\Adapt$ & Active plan; selected plan for opportunity $k$; initial/fallback adaptive plan.\\
$\LLM$ & Asynchronous LLM selector over the exact plans in $\Reg_k$.\\
$\sigma=(h,\eta,t_h)$ & Current step index, entry flag, and start time of its current dwell interval.\\
$J(o_k)$; $R(o_k)$ & Eligible target slots; observed contact anchors.\\
$\mathcal C$; $\mathcal A_{\rm rec}$ & Resource templates; recovery actions $\{\mathrm{spread},\mathrm{merge}\}$.\\
$T$; $\varepsilon$; $d_h$ & Soft total wall-time budget; native epoch cap; planned dwell of step $h$.\\
$\ell$; $\Delta F_\ell$; $\Delta P_\ell$ & An installed interval; incumbent gain; population objective-sum change.\\
\bottomrule
\end{tabularx}
\end{center}
Slot zero and the current best slot are protected from directed reseeding.
The population size, dwell durations, decision cadence, and admissible
decision horizon are fixed protocol parameters; their experimental values
are specified separately from the algorithms.

\clearpage
\begin{algorithm}[H]
\caption{CIPHEUR C05: asynchronous finite-plan control}
\label{alg:main}
\Input{Scheduling instance $\mathcal I$; population size $m$; LLM selector
$\LLM$; total budget $T$; epoch cap $\varepsilon$; fixed decision protocol.}
\Output{A feasible historical incumbent $X^*$.}
Construct or load $G$ from $\mathcal I$\;
$X^{(0)}\gets$ feasible degree-weight seed on $G$\;
$\Pop\gets\operatorname{NativeInitialize}(G,X^{(0)},m)$;
$X^*\gets\arg\max_{X\in\Pop\cup\{X^{(0)}\}}F(X)$\;
$\Trace\gets\Blank$; $\Hist\gets\Blank$; $k\gets0$\;
$o\gets\operatorname{Observe}(G,\Pop,X^*,\Trace,\Hist)$\;
$\pi\gets\Adapt$; $\sigma\gets\operatorname{StartPlan}(\pi)$;
$\ell\gets\operatorname{OpenInterval}(\pi,o,\Pop,X^*)$\;
\While{the total budget remains}{
  \If{an asynchronous decision $(k,\widehat\pi_k)$ has arrived}{
    \If{$\operatorname{ValidChoice}(\widehat\pi_k,o_k,\Reg_k)$ and time remains}{
      $\Hist\gets\Hist\mathbin{\|}\operatorname{CloseInterval}(\ell,\Pop,X^*,\sigma)$\;
      $\pi\gets$ the exact descriptor of $\widehat\pi_k$ in $\Reg_k$\;
      $\sigma\gets\operatorname{StartPlan}(\pi)$;
      $\ell\gets\operatorname{OpenInterval}(\pi,o_k,\Pop,X^*)$\;
    }
    Mark the decision request resolved\;
  }
  \If{a decision opportunity is due, no request is outstanding, and the protocol permits a new request}{
    $k\gets k+1$\;
    $\Hist_k\gets\operatorname{VisibleHistory}(\Hist)$\;
    $o_k\gets\operatorname{Observe}(G,\Pop,X^*,\Trace,\Hist_k,\pi,\sigma,\ell)$\;
    $\Reg_k\gets\operatorname{BuildRegistry}(o_k)$\tcp*[r]{Algorithm~\ref{alg:registry}}
    Request $\widehat\pi_k\gets\LLM(o_k,\Reg_k,\Hist_k)$ asynchronously, without waiting\;
  }
  \If{the total budget has expired}{\Break}
  $(\Pop,X^*,\sigma,z,f)\gets
  \operatorname{AdvancePlan}(G,\Pop,X^*,\pi,\sigma,\varepsilon)$\tcp*[r]{Algorithm~\ref{alg:execute}}
  Append the returned native progress record $z$, if present, to $\Trace$\;
  \If{$f$ requests fallback}{
    $\Hist\gets\Hist\mathbin{\|}\operatorname{CloseInterval}(\ell,\Pop,X^*,\sigma)$\;
    $o\gets\operatorname{Observe}(G,\Pop,X^*,\Trace,\Hist)$\;
    $\pi\gets\Adapt$; $\sigma\gets\operatorname{StartPlan}(\pi)$;
    $\ell\gets\operatorname{OpenInterval}(\pi,o,\Pop,X^*)$\;
  }
}
Disregard any unresolved decision; close the final interval into $\Hist$\;
Check that $X^*$ is independent in $G$, recompute $F(X^*)$, and check
$F(X^*)\geq F(X^{(0)})$\;
\Return{$X^*$}\;
\end{algorithm}

\textbf{Selection and execution.}
The asynchronous selector receives detached evidence. The native state is
owned by the search process, which continues Algorithm~\ref{alg:execute}
while a decision is outstanding. An accepted choice identifies a plan already
present in its observation-bound registry. Validation checks the binding,
registered descriptor identity, permitted fields, and freshness. Invalid
choices leave the active plan unchanged. A new request is admitted only at
an eligible decision opportunity with enough remaining budget.

\textbf{Observation.}
$o_k$ summarizes objective progress, stagnation, population diversity and
overlap, remaining time, sampled resource witnesses, target-slot blocker
evidence, recent closed intervals, and the partially observed active interval.
An observation is frozen for one decision; later installation does not
replace its slot or contact identifiers with newly preferred ones.

\clearpage
\begin{algorithm}[H]
\caption{BuildRegistry: observation-bound finite plans}
\label{alg:registry}
\Input{Frozen observation $o_k$; single-operation catalogue
$\mathcal A_0$; resource templates $\mathcal C$; recovery actions
$\mathcal A_{\rm rec}$; fixed dwell parameters.}
\Output{Finite registry $\Reg_k$ bound to $o_k$.}
$\Reg_k\gets\Blank$\;
\ForEach{$a\in\mathcal A_0$}{
  Add the registered single-step plan $[(a,d_a)]$ with its exact identifier to $\Reg_k$\;
}
$R(o_k)\gets$ the bounded set of genuine observed contact anchors\;
$J(o_k)\gets\{j: j\text{ has target-slot evidence in }o_k,\ j\neq0,
j\neq j_{\Best}(o_k)\}$\;
\ForEach{$j\in J(o_k)$}{
  \ForEach{$r\in R(o_k)$ with target-specific blocker evidence for slot $j$}{
    \ForEach{$c\in\mathcal C$ with the required resource labels available}{
      \ForEach{$a\in\mathcal A_{\rm rec}$}{
        $\pi\gets[(\operatorname{DirectedReseed}(j,r,c),d_{\rm seed}),
                         (a,d_a)]$\;
        Bind the exact target $j$, anchor $r$, template $c$, steps, and registered identifier to $\pi$\;
        $\Reg_k\gets\Reg_k\cup\{\pi\}$\;
      }
    }
  }
}
Freeze each descriptor and its identity with observation $o_k$\;
\Return{$\Reg_k$}\;
\end{algorithm}

The single-operation catalogue is
\[
\mathcal A_0=\{\mathrm{exploit},\mathrm{spread},\mathrm{merge},\mathrm{relay},
\mathrm{antenna},\mathrm{satellite},\mathrm{kick},\mathrm{reseed},\mathrm{adaptive}\}.
\]
The resource templates are $\mathcal C=\{\mathrm{antenna},\mathrm{satellite},
\mathrm{mixed}\}$. A directed plan first enters a specified trajectory with
a rooted resource-priority reseed, then performs the registered recovery
action. Full resource and time metadata are checked again at execution.

\textbf{Bounded evidence.}
For an observed anchor $r$ and target schedule $X_j$, blockers are exactly
\[
B_j(r)=N_G(r)\cap X_j,\qquad
W_j(r)=\sum_{v\in B_j(r)}w_v.
\]
Counts and total weights use the complete blocker set; displayed identifiers,
anchor candidates, and reuse summaries may be bounded samples. The anchor
selection is a predeclared sampling rule, and $w_r-W_j(r)$ is a context
quantity, not an observed repair gain. A plan does not modify $G$ or $w$.

\textbf{Whole-interval ledger.}
Opening an interval stores its binding observation, plan, installation time
$t_{\ell}^{-}$, incumbent value, and population objective sum. Closing it at
$t_{\ell}^{+}$ records
\begin{align*}
\Delta F_\ell
 &=F(X^*(t_{\ell}^{+}))-F(X^*(t_{\ell}^{-})),\\
\Delta P_\ell
 &=\sum_jF(X_j(t_{\ell}^{+}))-\sum_jF(X_j(t_{\ell}^{-})),\qquad
\Delta t_\ell=t_{\ell}^{+}-t_{\ell}^{-}.
\end{align*}
The ledger also retains the observed steps and whether execution was partial.
Intervals close upon replacement, fallback installation, or termination;
the active interval remains right-censored in the next observation. These
records describe realized search, including the recovery following an entry
intervention. They are not immediate seed-only rewards. The LLM uses these
records as decision context; this algorithm contains no LLM parameter update.

\clearpage
\begin{algorithm}[H]
\caption{AdvancePlan: bounded native execution and recovery}
\label{alg:execute}
\Input{$G$; population $\Pop$; historical incumbent $X^*$; active plan $\pi$;
step state $\sigma=(h,\eta,t_h)$; epoch cap $\varepsilon$.}
\Output{Updated $(\Pop,X^*,\sigma)$; optional progress record $z$;
fallback indicator $f$.}
$z\gets\Blank$; $f\gets\mathrm{false}$\;
\If{the planned dwell of step $h$ has elapsed}{
  \eIf{$\pi$ has a subsequent step}{
    $h\gets h+1$; $\eta\gets\mathrm{true}$\;
  }{
    Keep the final step $h$; $\eta\gets\mathrm{false}$\;
  }
  Restart the dwell clock $t_h$\;
  \Return{$(\Pop,X^*,(h,\eta,t_h),z,f)$}\;
}
$b\gets\min\{\varepsilon,\text{remaining total budget},
                        \text{remaining step dwell}\}$\;
$a\gets$ the active operation in step $h$\;
\eIf{$a=\operatorname{DirectedReseed}(j,r,c)$ and its binding or required metadata is invalid}{
  $\mathrm{status}\gets\mathrm{rejected}$\;
}{
  \eIf{$a=\operatorname{DirectedReseed}(j,r,c)$, $\eta=\mathrm{true}$,
         and $j\in\{0,j_{\Best}(\Pop)\}$}{
    $\mathrm{status}\gets\mathrm{rejected}$\;
  }{
    \If{$a=\operatorname{DirectedReseed}(j,r,c)$ and $\eta=\mathrm{true}$}{
      Install the registered rooted priority template $c$\;
      Attempt one bounded reseed of the exact slot $j$ around root $r$\;
      Check the resulting slot against the original graph and integer objective\;
    }
    Apply any single-action entry intervention only when $\eta=\mathrm{true}$\;
    Execute the registered native search operation using the unconsumed part of $b$\;
    Retain the native historical incumbent $X^*$ and obtain operation status\;
  }
}
Record the operation's objective, population, diversity, and execution effects in $z$\;
$\eta\gets\mathrm{false}$\;
\If{$\mathrm{status}\neq\mathrm{ok}$}{
  \eIf{$\pi$ has a subsequent step}{
    $h\gets h+1$; $\eta\gets\mathrm{true}$; restart $t_h$\;
  }{
    $f\gets\mathrm{true}$\;
  }
}
\Return{$(\Pop,X^*,(h,\eta,t_h),z,f)$}\;
\end{algorithm}

\textbf{Execution semantics.}
The reseed entry and its feasibility check share the epoch budget with native
search. The directed step's continuation is round-robin search over the
population; only the entry reseed has a specified target slot. A target that
has become protected is rejected, with its bound identifiers unchanged.
After the last planned dwell expires, the last search operation persists,
without repeating its entry intervention, until replacement or termination.
Native execution keeps the original weights and conflict edges, feasible
population states, and the best solution encountered so far.

The budget is a soft end-to-end wall-time limit: preparation, observation,
native operations, and final verification are included. The algorithms
abstract transport and persistence details while retaining the order of
observation, asynchronous selection, plan installation, native advancement,
and interval closure in the released C05 solver.

\end{document}
